\chapter{مقدمه}
\label{ch:intro}

\section{مقدمه کلی}

گراف‌ها ابزاری قدرتمند برای مدلسازی روابط بین موجودیت‌ها در بسیاری از سیستم‌های واقعی هستند. از شبکه‌های اجتماعی و زیستی گرفته تا شبکه‌های مخابراتی، حمل‌ونقل و حسگر بی‌سیم، ساختار گرافی نقش کلیدی در تحلیل و درک این سیستم‌ها ایفا می‌کند \cite{alanko2010}.

با افزایش اندازه و پیچیدگی گراف‌های دنیای واقعی، روش‌های سنتی تحلیل داده‌ها دیگر پاسخگوی این حجم از اطلاعات نیستند و نیاز به الگوریتم‌های تخصصی گراف بیش از پیش احساس می‌شود. این الگوریتم‌ها با بهره‌گیری از ویژگی‌های ساختاری گراف، امکان تحلیل، بهینه‌سازی و تصمیم‌گیری مؤثر را فراهم می‌کنند \cite{bengio2021, bondy2008}.

\section{مسائل بهینه‌سازی روی گراف}

در بسیاری از کاربردهای عملی، هدف صرفاً تحلیل ساختار گراف نیست، بلکه یافتن راه‌حل‌های بهینه برای مسائل تعریف‌شده بر روی آن اهمیت ویژه‌ای دارد. مسائل بهینه‌سازی گراف به یکی از شاخه‌های بنیادی علوم کامپیوتر، تحقیق در عملیات و هوش مصنوعی تبدیل شده‌اند. در این مسائل، هدف معمولاً انتخاب یا سازمان‌دهی مجموعه‌ای از رأس‌ها یا یال‌ها به‌گونه‌ای است که ضمن رعایت مجموعه‌ای از محدودیت‌ها، یک معیار مشخص بهینه شود \cite{paschos2010}.

از جمله مسائل شناخته‌شده در این حوزه می‌توان به کوتاه‌ترین مسیر، درخت پوشای کمینه، رنگ‌آمیزی گراف، پوشش رأس، خوشه‌بندی و مجموعه مستقل اشاره کرد. این مسائل در کاربردهای متنوعی نظیر طراحی شبکه‌های ارتباطی، زمان‌بندی، تخصیص منابع، مسیریابی، امنیت شبکه و تحلیل سیستم‌های پیچیده به‌کار می‌روند.

در میان این مسائل، مسئله مجموعه احاطه‌گر\footnote{Dominating Set} به‌دلیل کاربردهای گسترده و پیچیدگی محاسباتی بالا، جایگاه ویژه‌ای دارد. این مسئله در حوزه‌هایی مانند طراحی شبکه‌های حسگر بی‌سیم، مکان‌یابی تسهیلات، پایش شبکه‌های ارتباطی، کنترل و نظارت بر زیرساخت‌ها و بسیاری از سیستم‌های توزیع‌شده کاربرد فراوانی دارد \cite{haynes1998fundamentals, ore1962theory}.

\section{معرفی مسئله مجموعه احاطه‌گر}

مسئله مجموعه احاطه‌گر ($DS$) یکی از مسائل پایه‌ای نظریه گراف است. فرض کنید $G=(V,E)$ یک گراف ساده و متناهی باشد. مجموعه $D \subseteq V$ را یک مجموعه احاطه‌گر می‌نامیم هرگاه هر رأس $v \in V$ یا عضوی از $D$ باشد یا حداقل با یکی از رأس‌های موجود در $D$ مجاور باشد \cite{haynes1998fundamentals}.

هدف مسئله مجموعه احاطه‌گر کمینه، یافتن کوچک‌ترین مجموعه احاطه‌گر از نظر تعداد اعضا در گراف داده‌شده است. این مسئله برای حالت کلی NP-کامل است \cite{garey1979computers}.

\section{مسئله مجموعه احاطه‌گر کامل}

یکی از تعمیم‌های مهم مسئله مجموعه احاطه‌گر، مسئله مجموعه احاطه‌گر کامل\footnote{Total Dominating Set} ($TDS$) است. در این نسخه، شرط سخت‌گیرانه‌تر این است که هر رأس باید حداقل یک همسایه در مجموعه انتخاب‌شده داشته باشد؛ به‌بیان دیگر، هیچ رأسی نباید تنها توسط خودش احاطه شود و همچنین رأس‌های داخل مجموعه نیز باید توسط رأس دیگری از مجموعه احاطه گردند \cite{haynes1998fundamentals}.

فرموله‌سازی دقیق‌تر این مسئله در فصل \ref{ch:tdsp-extension} آمده است. مسئله مجموعه احاطه‌گر کامل نیز NP-کامل بوده و در بسیاری از کاربردهای عملی نظیر شبکه‌های حسگر بی‌سیم، کنترل ترافیک و سیستم‌های نظارتی اهمیت دارد.

\section{انگیزه پژوهش}

روش‌های دقیق برای حل مسائل NP-کامل تنها برای گراف‌های کوچک کارآمدند. با بزرگ شدن اندازه گراف‌ها، نیاز به روش‌های تقریبی و هوشمند افزایش می‌یابد. در سال‌های اخیر، ترکیب یادگیری عمیق با یادگیری تقویتی و شبکه‌های عصبی گرافی\footnote{Graph Neural Networks} توانسته است نتایج امیدوارکننده‌ای در حل مسائل ترکیبیاتی روی گراف ارائه دهد \cite{gilmer2017mpnn, khalil2017learning, bello2016neural}.

مقاله \cite{chen2024dsp} رویکردی مبتنی بر شبکه عصبی گرافی و یادگیری تقویتی دوگانه ($DDQN$) را برای مسئله مجموعه احاطه‌گر ارائه داده است. این پژوهش با هدف بررسی امکان‌پذیری گسترش همین چارچوب به مسئله مجموعه احاطه‌گر کامل و بهبود عملکرد آن انجام شده است.

\section{ساختار پایان‌نامه}

این پایان‌نامه در شش فصل تنظیم شده است. در فصل \ref{ch:theory} مبانی نظری لازم شامل نظریه گراف، یادگیری تقویتی و شبکه‌های عصبی گرافی بیان می‌شود. فصل \ref{ch:literature} به مرور پیشینه پژوهش می‌پردازد. فصل \ref{ch:reference} به بازسازی و پیاده‌سازی روش مرجع می‌پردازد. فصل \ref{ch:tdsp-extension} گسترش روش به مسئله مجموعه احاطه‌گر کامل و تحلیل نتایج را ارائه می‌دهد. فصل \ref{ch:conclusion} نیز شامل جمع‌بندی، دستاوردها، محدودیت‌ها و پیشنهادهای آینده است.
